\documentclass[11pt]{amsart}

\usepackage{amsmath,amssymb,amsthm}
\usepackage[margin=1in]{geometry}
\usepackage{hyperref}
\usepackage{mathtools}

\theoremstyle{plain}
\newtheorem{theorem}{Theorem}[section]
\newtheorem{proposition}[theorem]{Proposition}
\newtheorem{lemma}[theorem]{Lemma}
\newtheorem{corollary}[theorem]{Corollary}
\newtheorem{conjecture}[theorem]{Conjecture}

\theoremstyle{definition}
\newtheorem{definition}[theorem]{Definition}
\newtheorem{example}[theorem]{Example}
\newtheorem{remark}[theorem]{Remark}

\newtheoremstyle{gap}{}{}{\itshape}{}{\bfseries}{.}{ }{}
\theoremstyle{gap}
\newtheorem{gap}[theorem]{Gap (TODO for a prove session)}

\newcommand{\Poly}{\mathbf{Poly}}
\newcommand{\Cat}{\mathbf{Cat}}
\newcommand{\Cathash}{\mathbf{Cat}^{\sharp}}
\newcommand{\Set}{\mathbf{Set}}
\newcommand{\Comod}{\mathrm{Comod}}
\newcommand{\PRA}{\mathrm{PRA}}
\newcommand{\CLP}{\mathrm{CLP}}
\newcommand{\Fam}{\mathrm{Fam}}
\renewcommand{\lhd}{\vartriangleright}
\newcommand{\Store}[1]{\Delta #1}
\newcommand{\op}{^{\mathrm{op}}}
\newcommand{\yon}{\mathsf{y}}
\DeclareMathOperator{\const}{const}
\DeclareMathOperator{\ev}{ev}
\DeclareMathOperator{\im}{im}

\title[Workers, bicomodules, and parametric right adjoints]
{Workers, Bicomodules, and Parametric Right Adjoints:\\ a correspondence and a conjecture}

\author{MacBeth}
\address{Kodamai}
\date{20 September 2026}

\begin{document}

\begin{abstract}
Spivak, Garner and Fairbanks proved that bicomodules of polynomial comonads are
exactly the parametric right adjoints between the associated copresheaf categories,
and that every such functor factors as a connected-limit-preserving (right adjoint)
part followed by a coproduct (left adjoint) tail that carries no bicomodule data.
Separately, we have proved that the Workers state-composition discipline --- the
Dirichlet grading $\Store{S}\otimes\Store{T}=\Store{(S\times T)}$ on store comonads
--- is a canonical \emph{retract}, not a fibre, of the composition grading
$\Store{S}\lhd\Store{T}$ used by Braithwaite--Hedges--Mihejevs, the idempotent
collapsing each branch map to its self-evaluation. This note sets the two pictures
side by side. The retract's ``keep versus discard'' split has the same shape as the
factorization's ``right-adjoint content versus $\Sigma$-tail,'' which suggests the
identification \emph{worker $=$ bicomodule/parametric right adjoint}. We formulate
this as a precise conjecture and, rather than assert it, we isolate one concrete
obstacle --- an apparent \emph{reversal of adjointness} between the two splits ---
that a proof must confront. The correspondence is established literature; the
identification is a hypothesis under test.
\end{abstract}

\maketitle

\section{Introduction}

Two research programs describe composition of stateful, interacting systems in the
language of polynomial functors, and they have never been made to talk to each
other. This note is an attempt to build the bridge --- or, more honestly, to survey
the ground it must cross and to plant a flag where the hard step is.

\subsection*{The first picture: aggregation and parametric right adjoints}
In the polynomial-functor ecosystem \cite{NiuSpivak}, a small category is the same
thing as a polynomial comonad, and the natural maps \emph{between} such categories
come in two flavours: the ``vertical'' maps are retrofunctors (cofunctors), and the
``horizontal'' maps are two-sided \emph{bicomodules}. Spivak, Garner and Fairbanks
\cite{SGF} identify these bicomodules concretely: a bicomodule from the comonad $c$
to the comonad $d$ is exactly a \emph{parametric right adjoint} (equivalently a
connected-limit-preserving functor, equivalently a data-migration / aggregation
functor) between the copresheaf categories $c\text{-}\Set \to d\text{-}\Set$. Their
Proposition~3.6 factors any such functor into a genuinely ``polynomial'' right-adjoint
part and a coproduct tail that is not bicomodule data --- the aggregation ($\Sigma$)
step. This is the setting in which ``how does data flow and combine along a functor''
receives a clean categorical answer.

\subsection*{The second picture: workers and state}
In work aimed at the semantics of stateful agents \cite{workersProof}, one models a
piece of local state $S$ by the \emph{store comonad}, carried by the polynomial
$\Store{S}=S\cdot \yon^{S}$. Two workers with states $S$ and $T$ can be composed in
(at least) two ways. The \emph{Workers} discipline entangles their states via the
Dirichlet product, $\Store{S}\otimes\Store{T}=\Store{(S\times T)}$; the
Braithwaite--Hedges--Mihejevs (BHM) discipline nests them via functor composition,
$\Store{S}\lhd\Store{T}$, a store-of-stores. These are genuinely different
polynomials --- at $|S|=|T|=2$ they are $4\yon^{4}$ and $8\yon^{4}$ --- and our main
structural result about them is that the first is a canonical \emph{retract} of the
second (Theorem~\ref{thm:retract}, proved in \cite{workersProof}): there are
polynomial maps $\sigma$ and $r$ with $r\circ\sigma=\mathrm{id}$, and the idempotent
$e=\sigma\circ r$ collapses each branch map $g\colon S\to T$ to the constant at its
self-evaluation $g(s)$. The discarded data --- the non-constant branch maps --- is
exactly the store-of-store structure that nesting sees and entanglement forgets.

\subsection*{The question}
Both pictures split a rich construction into ``content that is kept'' and ``a tail
that is thrown away.'' In the aggregation picture the kept part is the right adjoint
and the tail is the coproduct $\Sigma$; in the workers picture the kept part is the
entangled state $\Store{(S\times T)}$ and the tail is the non-constant branch maps.
The shapes match. It is tempting --- and it has been repeatedly tempting, across
several reading sessions --- to conclude that the workers/BHM distinction is
\emph{precisely} the parametric-right-adjoint distinction: that a worker is a
bicomodule and the retract's kernel is the $\Sigma$-aggregation tail.

The purpose of this note is to state that temptation precisely and then to resist
asserting it. We recall the SGF correspondence (\S\ref{sec:corr}) and the workers
retract (\S\ref{sec:workers}) in a common language, formulate the identification as
Conjecture~\ref{conj:worker-pra} (\S\ref{sec:conj}), and --- this is the point ---
exhibit a specific obstacle to it. The obstacle is an \emph{apparent reversal of
adjointness}: the data the retract discards (branch maps $S\to T$) is
\emph{exponential}, hence right-adjoint-flavoured, whereas the data the SGF
factorization discards is a \emph{coproduct}, hence left-adjoint-flavoured. If the
identification is to hold, this reversal must be explained, not ignored. We record it
as a gap for a subsequent proof session, together with a second, type-level gap about
\emph{which} two comonads $(c,d)$ the workers construction is a bicomodule between.

\subsection*{Status of the claims}
The SGF correspondence is established, peer-reviewed literature
\cite{SGF}, deep-read for this note. The workers retract is our own result, proved
and machine-checked at small size \cite{workersProof}. The identification of the two
is a \emph{conjecture}; every statement in \S\ref{sec:conj} that goes beyond the two
established inputs is flagged as such. No new theorem is proved here.

\subsection*{Outline}
Section~\ref{sec:prelim} fixes the polynomial-functor conventions, the store comonad,
and the two products $\otimes,\lhd$. Section~\ref{sec:corr} states the SGF
correspondence and its factorization in our own words. Section~\ref{sec:workers}
recalls the workers retract. Section~\ref{sec:conj} states the conjecture and its two
gaps. Section~\ref{sec:conc} says what a proof session must settle.

\section{Preliminaries}\label{sec:prelim}

We work in the category $\Poly$ of polynomial functors $\Set\to\Set$, following the
conventions of \cite{NiuSpivak}.

\begin{definition}[polynomials and their maps]
A \emph{polynomial} is a coproduct of representables
$p=\sum_{s\in S_p}\yon^{p[s]}$, acting on sets by $p(X)=\sum_{s\in S_p}X^{p[s]}$; we
call $S_p$ its \emph{shapes} (or positions) and $p[s]$ its set of \emph{directions}
at $s$. A morphism $f\colon p\to q$ is a pair $f=(f_1,f^{\sharp})$ with
$f_1\colon S_p\to S_q$ on shapes and $f^{\sharp}_s\colon q[f_1 s]\to p[s]$
\emph{backward} on directions. Composition is $(g\circ f)_1=g_1 f_1$ and
$(g\circ f)^{\sharp}_s=f^{\sharp}_s\circ g^{\sharp}_{f_1 s}$: backward maps compose
in reverse. We track this reversal explicitly; it is where naïve shape-only reasoning
fails.
\end{definition}

$\Poly$ carries two monoidal structures we need.

\begin{definition}[the two products]
The \emph{composition product} is $(p\lhd q)(X)=p(q(X))$, with unit $\yon$; it is
strictly associative. The \emph{Dirichlet product} $p\otimes q$ has shapes
$S_p\times S_q$ and directions $p[s]\times q[t]$ at $(s,t)$, with unit $\yon$.
\end{definition}

\begin{definition}[polynomial comonads are categories]
A comonoid for $(\Poly,\lhd,\yon)$ is the same data as a small category
\cite{AhmanUustalu,NiuSpivak}: shapes are objects, directions at $s$ are the
morphisms out of $s$, the counit picks out identities and the comultiplication
composes. Under this dictionary the framed bicategory of polynomial comonads is
written $\Cathash$ \cite{NiuSpivak}. A polynomial comonad $c$ presents a category we
also call $c$, with copresheaf category $c\text{-}\Set$.
\end{definition}

\begin{definition}[store comonad]\label{def:store}
For a set $S$, the \emph{store} (or costate) polynomial is $\Store{S}:=S\cdot\yon^{S}$
--- shapes $S$, with $\Store{S}[s]=S$ at every shape. It carries the store comonad
$\Store{S}(X)=S\times X^{S}$, whose comultiplication is
$\delta_1(s)=(s,\mathrm{id}_S)$, $\delta^{\sharp}_s(i,j)=j$, and whose counit is
$\varepsilon_1(s)=\ast$, $\varepsilon^{\sharp}_s(\ast)=s$. As a category $\Store{S}$
is the \emph{codiscrete} (indiscrete) category on $S$: exactly one morphism between
any two objects.
\end{definition}

Finally, the object of study on the SGF side.

\begin{definition}[parametric right adjoint]\label{def:pra}
Let $C$ have a terminal object $1$. A functor $F\colon C\to D$ is a
\emph{parametric right adjoint} (PRA), or \emph{prafunctor}, if the first factor of
the canonical factorization
\[
C \xrightarrow{\;F/1\;} D/F1 \xrightarrow{\;U\;} D
\]
is a right adjoint, where $F/1$ sends $c$ to $Fc\to F1$ and $U$ forgets down to $D$
\cite[Def.~3.5]{SGF}. For endofunctors of $\Set$ this is equivalent to preserving
connected limits, and to being polynomial \cite[Prop.~2.5]{SGF}. A \emph{bicomodule}
from a polynomial comonad $c$ to a polynomial comonad $d$ is a polynomial equipped
with compatible left $c$- and right $d$-coactions; it is the ``horizontal'' notion of
map between the categories $c$ and $d$.
\end{definition}

\section{The Spivak--Garner--Fairbanks correspondence}\label{sec:corr}

We restate the two results of \cite{SGF} that this note leans on, in our own words.
Both are established; we claim no part of them.

\begin{theorem}[{\cite[Thm.~4.28]{SGF}}]\label{thm:sgf-main}
The framed bicategories $\Cathash$ and $\Comod(\Poly)$ are equivalent. Under the
equivalence: objects are small categories, i.e.\ polynomial comonads; the horizontal
(loose) maps from $c$ to $d$ are the $(c,d)$-bicomodules, which are exactly the
prafunctors $c\text{-}\Set\to d\text{-}\Set$; the vertical (tight) maps are the
comonoid homomorphisms, i.e.\ retrofunctors (not ordinary functors); and the $2$-cells
are natural transformations.
\end{theorem}

The identification of the horizontal maps is made pointwise by a chain of six
equivalent descriptions.

\begin{proposition}[{\cite[Prop.~3.9]{SGF}}]\label{prop:sixway}
For polynomial comonads $c,d$ there are equivalences
\[
\Comod(\Poly)(c,d)\;\simeq\;\PRA(c\text{-}\Set,\,d\text{-}\Set)
\;=\;\CLP(c\text{-}\Set,\,d\text{-}\Set)
\;\simeq\;[\,d,\ \Set[c]\,]
\;\simeq\;[\,d,\ \Fam((c\text{-}\Set)\op)\,],
\]
where $\CLP$ denotes the connected-limit-preserving functors.
\end{proposition}

The result we actually want to compare against workers is the factorization of a
single prafunctor, because it is what performs the ``keep versus discard'' split.

\begin{proposition}[{\cite[Prop.~3.6]{SGF}}, aggregation factorization]\label{prop:factor}
Every prafunctor $F\colon c\text{-}\Set\to\Set$ factors as
\[
c\text{-}\Set \xrightarrow{\;F/1\;} \Set^{F1} \xrightarrow{\;\Sigma\;} \Set,
\]
where $F/1$ is a \emph{right adjoint} --- the genuine polynomial/bicomodule content
--- and $\Sigma$ is the coproduct, a \emph{left adjoint} ``aggregation tail'' that
carries no bicomodule data.
\end{proposition}

The moral of Proposition~\ref{prop:factor} is exactly a keep/discard split: the
bicomodule content lives in the right-adjoint factor $F/1$, and the coproduct
$\Sigma$ is a tail one adds on top and which the correspondence does not see.

\begin{remark}[handedness]\label{rem:handedness}
Variance conventions differ across the literature: \cite{SGF} writes bicomodules as
$c\text{-}\Set\to d\text{-}\Set$, while Spivak's ``Categories by Kan extension''
writes the reversed $c'\text{-}\Set\to c\text{-}\Set$ via the opposite handedness on
bicomodules. It is the same theorem, but any concrete identification must first pin
the convention. We use the \cite{SGF} direction throughout.
\end{remark}

\section{The Workers retract}\label{sec:workers}

We recall the result of \cite{workersProof}. It is proved there in full and verified
computationally for shape sets of size $\le 3$; we import it and reuse its notation.

Write $A:=\Store{S}\otimes\Store{T}=\Store{(S\times T)}$ (the Workers grading) and
$B:=\Store{S}\lhd\Store{T}$ (the BHM grading). By Definition~\ref{def:store} and the
two product formulas, $A$ has shapes $S\times T$ with directions $S\times T$
everywhere, while $B$ has shapes $\{(s,g):s\in S,\ g\colon S\to T\}=S\times T^{S}$ and
directions $S\times T$ everywhere. So a shape of $B$ carries a whole \emph{branch
map} $g\colon S\to T$, where a shape of $A$ carries only a point $t\in T$.

\begin{theorem}[{\cite[Thm.~1.1]{workersProof}}]\label{thm:retract}
Define $\sigma\colon A\to B$ by $\sigma_1(s,t)=(s,\const_t)$ with identity backward
map, and $r\colon B\to A$ by $r_1(s,g)=(s,g(s))$ (\emph{self-evaluation}) with
identity backward map. Then $r\circ\sigma=\mathrm{id}_A$, while
$\sigma\circ r\ne\mathrm{id}_B$ whenever $|S|,|T|\ge 2$. Hence
$A=\Store{S}\otimes\Store{T}$ is a non-trivial retract of $B=\Store{S}\lhd\Store{T}$.
\end{theorem}

\begin{corollary}[the discarded data]\label{cor:idem}
The idempotent $e:=\sigma\circ r$ has $e_1(s,g)=(s,\const_{g(s)})$ and identity
backward map. It collapses each branch map $g$ to the constant at its self-evaluation
$g(s)$; its image is $A$. The data it discards is precisely the non-constant branch
maps $g\colon S\to T$ --- the store-of-store structure that $\lhd$ carries and
$\otimes$ forgets.
\end{corollary}

For orientation we also recall how the store comonad's own comultiplication sits over
the retraction; this is the sense in which ``$\otimes$ is the diagonal of $\lhd$.''

\begin{proposition}[{\cite[Thm.~3.1]{workersProof}}]\label{prop:diag}
Let $d\colon S\to S\times S$ be the diagonal and $\Store{d}\colon\Store{S}\to
\Store{(S\times S)}$ its store image (shape $s\mapsto(s,s)$, backward $\pi_2$). Then
$r_{S,S}\circ\delta=\Store{d}$: the store comultiplication $\delta$ is a lift of the
$\otimes$-diagonal along the retraction $r$.
\end{proposition}

Corollary~\ref{cor:idem} is the crux for the comparison. The workers retract keeps a
point-valued datum ($g(s)\in T$, equivalently the entangled state $S\times T$) and
discards the branch map's dependence on inputs other than the current one.

\section{The conjecture, and where it breaks}\label{sec:conj}

We can now line the two keep/discard splits up.

\begin{center}
\begin{tabular}{lll}
\hline
 & \textbf{kept} & \textbf{discarded (tail)}\\
\hline
SGF factorization (Prop.~\ref{prop:factor}) & $F/1$: right adjoint & $\Sigma$: coproduct (left adjoint)\\
Workers retract (Cor.~\ref{cor:idem}) & $A=\Store{(S\times T)}$ & non-constant branch maps $g\colon S\to T$\\
\hline
\end{tabular}
\end{center}

The parallel is real: in each row a rich construction is split into a piece the
relevant correspondence recognises and a tail it does not. It motivates the following.

\begin{conjecture}[worker $=$ bicomodule/PRA]\label{conj:worker-pra}
There is a choice of polynomial comonads $(c,d)$ for which the BHM grading
$B=\Store{S}\lhd\Store{T}$ presents a $(c,d)$-bicomodule, equivalently a prafunctor
$c\text{-}\Set\to d\text{-}\Set$; under the aggregation factorization of
Proposition~\ref{prop:factor}, the retract $A=\Store{S}\otimes\Store{T}$ of
Theorem~\ref{thm:retract} is the right-adjoint (bicomodule-content) factor, and the
data discarded by the idempotent $e$ of Corollary~\ref{cor:idem} is the coproduct
tail $\Sigma$. Consequently ``worker'' $=$ ``parametric right adjoint / bicomodule,''
and the retract kernel $=$ the aggregation tail.
\end{conjecture}

We do not claim Conjecture~\ref{conj:worker-pra}. Two gaps stand between the suggestive
table and a proof; naming them precisely is the real content of this section.

\begin{gap}[the apparent reversal of adjointness]\label{gap:adjoint}
The two ``discarded tails'' point in opposite adjoint directions, and the conjecture
as stated papers over this. In Proposition~\ref{prop:factor} the discarded factor is
$\Sigma$, a \emph{coproduct}, hence a \emph{left} adjoint; the \emph{kept} factor
$F/1$ is the right adjoint. In Corollary~\ref{cor:idem} the discarded data is the set
of branch maps $g\colon S\to T$. But branch maps are \emph{exponential} data: they
are the elements of $T^{S}$, and the power $(-)^{S}$ is a \emph{right} adjoint (right
adjoint to the copower $S\times(-)$) --- the adjoint flavour the SGF factorization
\emph{keeps} in its $F/1$ factor. On the face of it, then, the workers retract discards right-adjoint
content while the SGF factorization discards left-adjoint content. Either (i) the
correct comparison reverses one of the two splits (e.g.\ via the handedness of
Remark~\ref{rem:handedness}, or by comparing $\sigma$ rather than $e$), so that the
exponential branch data reappears on the kept side; or (ii) the analogy is only
formal and Conjecture~\ref{conj:worker-pra} is false. Deciding which is the first
task of a proof session. A concrete sub-question: is the exponential $T^S$ in a
$B$-shape the connected-limit content that $F/1$ retains, with the genuine
$\Sigma$-tail located elsewhere (e.g.\ in the coproduct $S\cdot(-)$ that makes
$\Store{S}=S\cdot\yon^S$ a sum of representables in the first place)?
\end{gap}

\begin{gap}[which two comonads?]\label{gap:cd}
Proposition~\ref{prop:sixway} is about maps \emph{between fixed comonads},
$\Comod(\Poly)(c,d)$, whereas $\Store{S}\otimes\Store{T}$ and
$\Store{S}\lhd\Store{T}$ are \emph{monoidal products of two comonads}, producing new
polynomials rather than a map between two others. Before the correspondence can bite,
one must say what plays the role of $(c,d)$. Candidates: (a) take $c=\Store{T}$,
$d=\Store{S}$ and read $\Store{S}\lhd\Store{T}$ as (the carrier of) a bicomodule via
the composite comonad structure --- but $\lhd$ of two comonads need not be a comonad,
so this requires care; (b) fix $c=d=\yon$ (the terminal category) and read everything
as an endo-prafunctor of $\Set$, in which case the correspondence degenerates and the
comparison is only with the single-variable Proposition~\ref{prop:factor}; (c) use
the store comonads' status as codiscrete categories (Definition~\ref{def:store}) to
make $\Store{S},\Store{T}$ the two objects and the graded composition a genuine
bicomodule between them. Only (c) has a chance of engaging the full
Theorem~\ref{thm:sgf-main}; verifying that the BHM grading \emph{is} the
corresponding bicomodule is not done here.
\end{gap}

\begin{remark}[the nearest neighbour, and why it is not a shortcut]
The graded-comonad reading of workers has an established relative:
Capucci--Myers's ``contextads'' \cite{ContextadsWreaths}, where a polynomial monad
transposes to a dependently graded comonad and composition multiplies the grade. Its
Theorem~A.4 ($\mathrm{Kl}(T)\simeq\mathrm{Ctx}(\odot)$ for a parametric right adjoint
monad $T$) is stated for exactly the PRA class that appears here --- but it is left
\emph{unproven} in \cite{ContextadsWreaths}. So the literature already gestures at
``graded composition lives in the PRA world,'' but it does not supply the
bicomodule identification we need, and cannot be cited as a proof of
Conjecture~\ref{conj:worker-pra}. If anything, settling Gaps~\ref{gap:adjoint}
and~\ref{gap:cd} would bear on Theorem~A.4, not the other way round.
\end{remark}

\section{What a proof session must settle}\label{sec:conc}

The correspondence of Spivak--Garner--Fairbanks and the workers retract are each
solid; the bridge between them is not yet built. This note has reduced ``is a worker a
bicomodule?'' to two precise tasks.

\begin{enumerate}
\item[(1)] \textbf{Resolve the adjointness reversal (Gap~\ref{gap:adjoint}).} The
exponential branch data $T^{S}$ that the retract discards is right-adjoint-flavoured,
whereas the SGF factorization discards a coproduct. Determine whether a handedness or
$\sigma$-versus-$e$ change of comparison puts the two splits on the same side, or
whether the analogy fails.
\item[(2)] \textbf{Pin the comonads (Gap~\ref{gap:cd}).} Exhibit the pair $(c,d)$ for
which $\Store{S}\lhd\Store{T}$ is a $(c,d)$-bicomodule, presumably using the store
comonads as codiscrete categories, and check that the corresponding prafunctor's
Proposition~\ref{prop:factor} factorization has the workers retract as its
right-adjoint factor.
\end{enumerate}

If both succeed, Conjecture~\ref{conj:worker-pra} holds and ``worker'' acquires a
clean semantics inside $\Poly$: a worker is a parametric right adjoint, its
composition is bicomodule composition, and the retract that separates entanglement
from nesting is the aggregation factorization. That is worth having --- it would place
the state axis of the compositional-correctness program inside the same framed
bicategory $\Cathash$ that already houses the directed (retrofunctor) axis. If either
fails, the failure is itself informative: it would say that state-entanglement is a
\emph{different} kind of collapse from data-aggregation, and would sharpen the
question of what the store-of-store structure really is. Both outcomes are progress;
neither is available without the proof session.

\begin{thebibliography}{9}

\bibitem{AhmanUustalu}
D.~Ahman and T.~Uustalu,
\emph{Directed containers as categories},
in \emph{Proc.\ MSFP 2016}, EPTCS 207, 2016, pp.~89--98.
arXiv:1604.01187.

\bibitem{ContextadsWreaths}
M.~Capucci and D.~J.~Myers,
\emph{Contextads as wreaths; Kleisli, Para, and Span constructions},
arXiv:2410.21889, 2024.

\bibitem{NiuSpivak}
N.~Niu and D.~I.~Spivak,
\emph{Polynomial Functors: A Mathematical Theory of Interaction},
arXiv:2312.00990, 2023.

\bibitem{SGF}
D.~I.~Spivak, R.~Garner and A.~Fairbanks,
\emph{Functorial aggregation},
J.\ Pure Appl.\ Algebra \textbf{229} (2025), no.~2, art.~107883.
arXiv:2111.10968.

\bibitem{workersProof}
MacBeth,
\emph{The Workers $\otimes$-grading is a retract of the BHM $\lhd$-grading},
Kodamai internal proof note, 2026-08-29.
Prior art: J.~Braithwaite, J.~Hedges and D.~Mihejevs (ACT 2026), graded monad
$T_P(X)=X\lhd P$.

\end{thebibliography}

\end{document}
